\documentclass[11pt]{article}
\usepackage{amsmath,amssymb,amsfonts}
\usepackage{physics}
\usepackage{geometry}
\usepackage{hyperref}
\usepackage{cite}
\usepackage{bm}
\title{Measurement, Operators, and Incompatibility}
\begin{document}
\maketitle

\section{Primitive Structure: Events, Not Operators}

\begin{axiom}[Quantum event structure]
A quantum measurement is specified by a normalized POVM
\[
E:\mathcal B(\Omega)\to\mathcal B(\mathcal H),
\]
where $\Omega$ is the outcome space.
\end{axiom}

\begin{axiom}[Born rule]
For any state $\rho$ and event $\Delta\subset\Omega$,
\[
\mathbb P(\Delta|\rho)=\mathrm{Tr}(\rho E(\Delta)).
\]
\end{axiom}

\begin{remark}
No linear operator by itself represents a measurement.
Measurements are event-valued, not operator-valued.
\end{remark}

---

\section{Operators as Derived Objects}

\begin{definition}[Observable operator]
Given a PVM $E_A$, define the associated operator
\[
\hat A := \int_\Omega a\, dE_A(a).
\]
\end{definition}

\begin{proposition}
The operator $\hat A$ encodes only the moments of the probability measure
$\mu_\rho^A(\Delta)=\mathrm{Tr}(\rho E_A(\Delta))$.
\end{proposition}

\begin{corollary}
Operator action $\hat A\psi$ has no direct operational meaning as a
measurement outcome.
\end{corollary}

\begin{center}
\emph{Measurement $\Rightarrow$ PVM/POVM $\Rightarrow$ operator, never the reverse.}
\end{center}

---

\section{Definiteness and Eigenstates}

\begin{definition}[Definite value]
A state $\psi$ has a definite value of observable $A$ iff
\[
E_A(\Delta)\psi=\psi
\]
for some measurable $\Delta$.
\end{definition}

\begin{theorem}[Continuous observables have no sharp states]
If $E_A$ has continuous spectrum, then
\[
E_A(\{a\})=0 \quad \forall a,
\]
hence no normalizable state has a definite value of $A$.
\end{theorem}

\begin{remark}
Generalized eigenstates (Dirac deltas, plane waves) define the PVM but are
not physical states.
This structure originates in the formalism of
\emph{:contentReference[oaicite:0]{index=0}}.
\end{remark}

---

\section{Why Multiplication Does Not Imply Definiteness}

For position on $\Omega$,
\[
(E_X(\Delta)\psi)(x)=\chi_\Delta(x)\psi(x),
\qquad
(\hat X\psi)(x)=x\psi(x).
\]

\begin{theorem}
A wavefunction $\psi$ is an eigenstate of $\hat X$ iff it is supported on a
single point, hence is distributional.
\end{theorem}

\begin{proof}
The eigenvalue equation $x\psi(x)=\lambda\psi(x)$ implies
$(x-\lambda)\psi(x)=0$ almost everywhere.
\end{proof}

\begin{corollary}
Multiplication operators do not endow states with definite values.
Definiteness is determined solely by spectral projections.
\end{corollary}

---

\section{Incompatibility and Noncommutativity}

\begin{definition}[Joint measurability]
Two observables $A,B$ are jointly sharply measurable iff there exists a
joint PVM $E_{AB}$ with correct marginals.
\end{definition}

\begin{theorem}[Incompatibility criterion]
$A$ and $B$ admit a joint PVM iff
\[
[E_A(\Delta),E_B(\Gamma)]=0
\quad \forall \Delta,\Gamma.
\]
\end{theorem}

\begin{corollary}
If $[\hat A,\hat B]\neq 0$, then no joint sharp measurement exists.
\end{corollary}

\begin{remark}
Operator noncommutativity is a diagnostic for measurement incompatibility,
not its definition.
\end{remark}

---

\section{Heisenberg Uncertainty Without Operators}

\begin{theorem}[Uncertainty as POVM obstruction]
There exists no POVM on phase space $\mathbb R^2$ whose marginals are both
arbitrarily sharp position and momentum measurements.
\end{theorem}

\begin{corollary}
Heisenberg uncertainty expresses the nonexistence of a joint refinement of
the position and momentum POVMs.
\end{corollary}

\begin{remark}
The usual inequality $\Delta X\,\Delta P\ge\hbar/2$ is a corollary of this
structural incompatibility, not its cause.
This insight clarifies the original intuition of
\emph{:contentReference[oaicite:1]{index=1}}.
\end{remark}

---

\section{POVMs and Finite Resolution}

\begin{definition}[Smeared measurement]
A realistic measurement is described by a POVM
\[
F(\Delta)=\int \mu(\Delta|x)\, dE(x),
\]
where $\mu$ models detector resolution.
\end{definition}

\begin{proposition}
POVMs allow approximate joint measurements even when sharp PVMs are
incompatible.
\end{proposition}

---

\section{Collapse as Conditioning, Not Dynamics}

\begin{axiom}[State update]
Upon outcome $\Delta$, the state updates as
\[
\rho \mapsto
\frac{\sqrt{E(\Delta)}\,\rho\,\sqrt{E(\Delta)}}
{\mathrm{Tr}(\rho E(\Delta))}.
\]
\end{axiom}

\begin{theorem}[No-collapse theorem]
State update under measurement is not generated by any Hamiltonian and is
not a physical time evolution.
\end{theorem}

\begin{corollary}
Collapse is Bayesian conditioning on a non-Boolean event structure.
\end{corollary}

---

\section{Bundle Interpretation}

\begin{proposition}
Position is a base-manifold observable.
Momentum arises from a connection on the state bundle.
\end{proposition}

\begin{corollary}
The asymmetry between position and momentum is geometric, not interpretive.
\end{corollary}

---

\section{Final Synthesis}

\begin{itemize}
\item Measurements are POVMs/PVMs, not operators.
\item Operators summarize measurements but do not define them.
\item Noncommutativity means incompatibility of event structures.
\item Uncertainty is the absence of a joint POVM.
\item Collapse is information update, not dynamics.
\end{itemize}

\begin{center}
\boxed{
\text{Quantum mechanics is probability theory on a non-Boolean event space.}
}
\end{center}
\end{document}
